```latex
\documentclass[11pt,a4paper]{article}

\usepackage[margin=1in]{geometry}
\usepackage{amsmath,amssymb,amsthm,mathtools}
\usepackage{enumitem}
\usepackage{hyperref}
\usepackage{microtype}

\title{A Zero-Preserving, Singularity-Sensitive Division Calculus with an Idempotent Element}

\author{Manu Iyengar}

\date{\today}

\newtheorem{definition}{Definition}[section]
\newtheorem{proposition}[definition]{Proposition}
\newtheorem{theorem}[definition]{Theorem}
\newtheorem{remark}[definition]{Remark}
\newtheorem{example}[definition]{Example}

\newcommand{\I}{I}
\newcommand{\Div}{\operatorname{Div}}

\begin{document}

\maketitle

\begin{abstract}

This paper proposes a division calculus designed to preserve ordinary
multiplication by zero while assigning a meaningful value to division by
zero. A distinguished element \(I\) is introduced with \(I^2=I\), ordinary
multiplication by zero remains \(0x=x0=0\), and division by zero is defined
by \(a/0=aI\), giving \(0/0=0\).

The central feature is that division is syntax-sensitive: denominator
factors are retained long enough for zero factors to be detected before
ordinary algebraic simplification collapses them. Consequently, for
\(c\neq0\),

\[
\frac{x}{0c}=\frac{xI}{c},
\]

rather than \(x/(0c)=x/0\).

The underlying algebra can be represented as
\(\mathbb{R}[I]/(I^2-I)\), which is isomorphic to
\(\mathbb{R}\times\mathbb{R}\). The resulting framework is therefore not
proposed as a conventional field-like algebra with a total binary division
operation. Instead, it is a structured division calculus over an ordinary
algebraic carrier.

The paper formalizes this distinction, gives a proposed syntax and
evaluation rule, identifies valid and invalid identities, and places the
proposal in relation to wheels, meadows, and related approaches to division
by zero.

\end{abstract}


\section{Introduction}

Division by zero is undefined in ordinary real arithmetic because the
equation \(bx=a\) has no unique solution when \(b=0\).

Numerous algebraic frameworks have nevertheless investigated ways of
extending arithmetic so that division becomes total or singular expressions
receive designated values. Examples include wheels
\cite{Carlstrom2004}, meadows
\cite{BergstraHirshfeldTucker2009}, non-involutive meadows
\cite{BergstraMiddelburg2015}, and common meadows
\cite{BergstraPonse2015}.

The present proposal starts from a different design requirement. We wish
to retain the ordinary rule

\[
0x=x0=0
\]

while also assigning a value to division by zero.

Let \(I\) denote the distinguished singular element associated with
\(1/0\). The basic proposed rules are

\[
I^2=I,
\qquad
0I=I0=0,
\qquad
\frac{a}{0}=aI.
\]

In particular,

\[
\frac{0}{0}=0.
\]

These rules alone are not the main difficulty. The crucial issue appears
when zero occurs as a factor inside a denominator.

If \(c\neq0\), the proposal is

\[
\frac{x}{0c}=\frac{xI}{c},
\]

whereas ordinary multiplication gives \(0c=0\).

Hence the expression \(0c\) cannot be collapsed to \(0\) before division
semantics are applied.

This observation forces a distinction between ordinary algebraic equality
and equality of structured division expressions.


\section{Motivation and Design Requirements}

The proposed calculus is motivated by four requirements.

\subsection{Division by zero receives a value}

For every algebraic element \(a\), the expression \(a/0\) should have a
defined result:

\[
\frac{a}{0}=aI.
\]

\subsection{Multiplication by zero remains ordinary}

The framework should not replace the usual absorbing behavior of zero:

\[
0x=x0=0.
\]

In particular,

\[
0I=I0=0.
\]

This requirement is deliberately different from treating \(I\) as an
ordinary multiplicative inverse of zero.

The notation \(I=1/0\) is a definition of the singular division value, not
an assertion that

\[
I\cdot0=1.
\]

\subsection{Ordinary nonzero division remains ordinary}

For \(b\neq0\),

\[
\frac{a}{b}
\]

has its usual meaning, and

\[
\left(\frac{a}{b}\right)b=a.
\]

\subsection{Zero factors in denominators remain visible}

If \(c\neq0\), the proposed semantics distinguishes

\[
\frac{x}{0}
\]

from

\[
\frac{x}{0c}.
\]

Specifically,

\[
\frac{x}{0}=xI,
\qquad
\frac{x}{0c}=\frac{xI}{c}.
\]


\section{The Underlying Algebra}

\subsection{Definition of the singular element}

Introduce a distinguished element \(I\) satisfying

\[
I^2=I.
\]

The symbol \(I\) is associated with division by zero by the definition

\[
\frac{1}{0}:=I.
\]

No inverse law is inferred from this definition.

\subsection{The algebraic carrier}

Consider

\[
A=\mathbb{R}[I]/(I^2-I).
\]

Every element of \(A\) has a representative of the form

\[
a+bI,
\qquad
a,b\in\mathbb{R}.
\]

Addition is

\[
(a+bI)+(c+dI)
=
(a+c)+(b+d)I,
\]

and multiplication is

\[
(a+bI)(c+dI)
=
ac+(ad+bc+bd)I.
\]

The ordinary absorbing property of zero follows immediately:

\[
0(a+bI)=0.
\]


\subsection{Representation as a product algebra}

There is an explicit isomorphism

\[
\phi:A\longrightarrow\mathbb{R}\times\mathbb{R}
\]

defined by

\[
\phi(a+bI)=(a,a+b).
\]

Indeed,

\[
\phi(I)=(0,1),
\]

and

\[
(0,1)^2=(0,1).
\]

The element \(I\) is therefore an idempotent and a zero divisor:

\[
I(1-I)=0.
\]

This is an important structural fact. The proposal does not claim that
the resulting carrier is a field or an integral domain.


\section{Why Division Must Be Syntax-Sensitive}

\subsection{The \(0=0c\) problem}

In the underlying algebra,

\[
0=0c
\]

for every \(c\).

If division were an ordinary extensional binary operation on algebraic
values, substitutivity would imply

\[
\frac{x}{0}=\frac{x}{0c}.
\]

The proposed semantics instead gives, for \(c\neq0\),

\[
\frac{x}{0}=xI
\]

and

\[
\frac{x}{0c}=\frac{xI}{c}.
\]

For example,

\[
\frac{1}{0}=I,
\qquad
\frac{1}{0\cdot2}=\frac{I}{2}.
\]

Thus unrestricted substitution of algebraically equal expressions inside
a pending division expression is incompatible with the intended semantics.


\subsection{Algebraic equality versus division-expression equivalence}

The framework therefore distinguishes two notions.

First, ordinary algebraic equality is used after an expression has been
evaluated into the carrier \(A\).

Second, structured division expressions retain denominator-factor
information until the division rule has been applied.

This is analogous to the distinction between a source expression and its
evaluated value in a programming-language semantics: two expressions can
have equal values in one context while still carrying different syntactic
information needed by a context-sensitive operation.

\begin{remark}

This distinction is not a claim that ordinary equality in \(A\) is
non-transitive or otherwise defective.

Rather, it says that algebraic equality of already-evaluated values cannot
be used as an unrestricted preprocessing rule for structured division
expressions.

\end{remark}


\section{Syntax of the Division Calculus}

\subsection{Ordinary algebraic expressions}

Let ordinary algebraic expressions be generated by

\[
E ::= r \mid I \mid (E+E) \mid (E\cdot E),
\qquad r\in\mathbb{R}.
\]

\subsection{Structured division expressions}

Instead of immediately representing a denominator as a single algebraic
value, a division expression retains its denominator factors:

\[
\Div(N;d_1,\ldots,d_n).
\]

Informally this is written as

\[
\frac{N}{d_1\cdots d_n}.
\]

The structured representation records the factors before singular
division semantics are applied.


\section{Division Semantics}

Let

\[
Z=\{i:d_i=0\}.
\]

For a numerator \(x\), define

\[
\Div(x;d_1,\ldots,d_n)
=
xI^{|Z|}
\prod_{i\notin Z}d_i^{-1},
\]

where the inverses in the final product are ordinary inverses of nonzero
denominators.

Since

\[
I^k=I
\qquad
\text{for every }k\geq1,
\]

this becomes

\[
\Div(x;d_1,\ldots,d_n)
=
\begin{cases}
x\displaystyle\prod_i d_i^{-1},
& Z=\varnothing,\\[1.2ex]
xI\displaystyle\prod_{d_i\neq0}d_i^{-1},
& Z\neq\varnothing.
\end{cases}
\]


\subsection{Division by a nonzero denominator}

For \(c\neq0\),

\[
\Div(x;c)=xc^{-1}.
\]

Thus ordinary division is unchanged whenever no zero denominator factor
is present.


\subsection{Division by zero}

For the single-factor case,

\[
\Div(x;0)=xI.
\]

Consequently,

\[
\frac{0}{0}=0I=0.
\]


\subsection{A zero factor followed by a nonzero factor}

For \(c\neq0\),

\[
\Div(x;0,c)=xIc^{-1}.
\]

In conventional notation,

\[
\frac{x}{0c}=\frac{xI}{c}.
\]


\subsection{Multiple zero factors}

For example,

\[
\Div(x;0,0,c)
=
xI^2c^{-1}
=
xIc^{-1}.
\]

More generally, the number of zero factors does not matter after the first
one because \(I\) is idempotent:

\[
I^n=I,
\qquad
n\geq1.
\]

Thus

\[
\frac{x}{0\cdot0\cdot c}
=
\frac{xI}{c}.
\]


\section{Algebraic Properties}

\subsection{Zero divided by zero}

The defining rule gives

\[
\frac{0}{0}=0I=0.
\]

This is one of the principal design choices of the framework.


\subsection{Multiplication after division}

For \(b\neq0\),

\[
\left(\frac{a}{b}\right)b=a.
\]

For division by zero,

\[
\left(\frac{a}{0}\right)0
=
(aI)0
=
0.
\]

Therefore the identity

\[
\left(\frac{a}{b}\right)b=a
\]

cannot be used without the condition \(b\neq0\).


\subsection{No cancellation through zero}

From

\[
\left(\frac{a}{0}\right)0=0
\]

one cannot recover \(a\).

Thus cancellation through zero is invalid.

Similarly, since \(I\) is a zero divisor,

\[
I(1-I)=0,
\]

cancellation involving \(I\) is not generally valid.


\subsection{Some valid singular identities}

The following identities follow from the proposed rules:

\[
\frac{a}{0}+\frac{b}{0}
=
(a+b)I
=
\frac{a+b}{0},
\]

and

\[
\frac{a}{0}\frac{b}{0}
=
abI^2
=
abI
=
\frac{ab}{0}.
\]

These identities illustrate that the singular element behaves coherently
under the ordinary operations of the carrier.


\subsection{Restricted fraction identities}

The usual identity

\[
\frac{a}{b}+\frac{c}{d}
=
\frac{ad+bc}{bd}
\]

cannot be assumed universally in the presence of structured singular
denominators.

In particular, combining denominator factors can destroy information
about which factors were zero.

Any such transformation therefore requires a proof that it preserves the
structured division semantics.


\section{Rewrite Discipline}

The proposed evaluation procedure is:

\begin{enumerate}[label=\arabic*.]

\item Parse the division expression while retaining denominator factors.

\item Do not collapse zero-containing denominator factors using ordinary
algebraic multiplication.

\item Detect which denominator factors are zero.

\item Replace each zero factor by its singular \(I\)-contribution.

\item Evaluate all remaining nonzero denominator factors using ordinary
inverses.

\item Normalize the resulting algebraic expression in \(A\).

\end{enumerate}

In particular, the rewrite

\[
0c\longrightarrow0
\]

is valid in ordinary algebra but is not permitted as a preprocessing step
inside a pending structured division expression.


\section{Valid and Invalid Transformations}

The following are valid:

\begin{align*}
I^2 &= I,\\
0I &= 0,\\
\frac{a}{0} &= aI,\\
\frac{0}{0} &= 0,\\
\left(\frac{a}{b}\right)b &= a
\qquad (b\neq0),\\
\left(\frac{a}{0}\right)0 &= 0,\\
\frac{x}{0c} &= \frac{xI}{c}
\qquad (c\neq0),\\
\frac{x}{0\cdot0\cdot c} &= \frac{xI}{c}
\qquad (c\neq0).
\end{align*}

The following are invalid or require explicit restrictions:

\begin{align*}
\left(\frac{a}{b}\right)b &= a
\qquad\text{without }b\neq0,\\
\left(\frac{a}{0}\right)0 &= a,\\
\frac{x}{0c} &= \frac{x}{0}
\qquad(c\neq0),\\
\text{cancellation through }0,\\
\text{cancellation through }I.
\end{align*}

The usual fraction-combination identity

\[
\frac{a}{b}+\frac{c}{d}
=
\frac{ad+bc}{bd}
\]

must also be treated as conditional rather than universally valid.


\section{Relation to Existing Frameworks}

\subsection{Wheels}

Wheels provide an algebraic approach to totalized division in which
division by zero receives designated behavior and ordinary field laws are
modified appropriately \cite{Carlstrom2004}.

The present proposal differs in its emphasis on retaining denominator
factor structure during evaluation.


\subsection{Meadows}

Meadows provide an equational specification of totalized division
\cite{BergstraHirshfeldTucker2009}.

Their semantics is designed so that division can be treated algebraically
within the chosen signature.

The proposed calculus instead deliberately restricts unrestricted
extensional treatment of division because denominator syntax carries
semantic information.


\subsection{Non-involutive and common meadows}

Further variants study different choices for the behavior of inverse and
division at zero \cite{BergstraMiddelburg2015,BergstraPonse2015}.

These provide important comparison points for determining exactly which
properties of the present proposal are genuinely distinct.


\subsection{Idempotent singular extensions}

Idempotent extensions involving a distinguished element associated with
\(1/0\) have also appeared in other mathematical proposals.

Therefore the mere introduction of an idempotent \(I\) is not claimed as
novel.

The potentially distinctive aspect investigated here is the combination
of:

\begin{enumerate}

\item ordinary zero absorption \(0I=0\);

\item the definition \(a/0=aI\);

\item the resulting \(0/0=0\);

\item and factor-sensitive, syntax-preserving division semantics.

\end{enumerate}

A complete novelty assessment requires a substantially broader literature
review than the preliminary comparison presented here.


\section{Novelty and Scope}

The proposal should be viewed conservatively as a research program rather
than as an established new algebraic theory.

The most important open question is whether the syntax-sensitive semantics
can be formalized in a way that is mathematically rigorous, useful, and
sufficiently distinct from existing systems.

In particular, a complete treatment should establish:

\begin{itemize}

\item a precise grammar for all expressions;

\item a formally defined evaluation function;

\item an appropriate equivalence relation;

\item termination of the proposed normalization procedure;

\item confluence or a characterization of unavoidable non-confluence;

\item treatment of nested divisions;

\item compatibility with substitution;

\item and a rigorous comparison with existing totalized-division
frameworks.

\end{itemize}


\section{Open Problems}

Several questions remain open.

\subsection{Nested division}

Expressions such as

\[
\frac{x}{(y/0)}
\]

require a precise decision about whether the denominator is analyzed as a
structured expression before or after the inner division is evaluated.


\subsection{Normal forms}

A canonical normal form for arbitrary expressions containing structured
division has not yet been established.


\subsection{Rewrite properties}

It remains to be proved whether the proposed rewrite discipline is
terminating and confluent.


\subsection{Substitution}

The precise conditions under which ordinary algebraic substitution is
allowed inside a structured division expression require formal treatment.


\subsection{General base rings}

The construction has been described over \(\mathbb{R}\), but much of the
underlying algebra can be formulated over a more general commutative ring.

Determining the minimal assumptions required for the division calculus is
an open direction.


\section{Conclusion}

This paper proposes a zero-preserving approach to division by zero based
on an idempotent singular element \(I\).

The carrier

\[
\mathbb{R}[I]/(I^2-I)
\]

provides a simple algebraic environment in which

\[
I^2=I,
\qquad
0I=0,
\]

while the division rule

\[
a/0=aI
\]

gives

\[
0/0=0.
\]

The central proposal is not merely the introduction of \(I\), but the
retention of denominator-factor structure during division.

This yields

\[
\frac{x}{0c}=\frac{xI}{c}
\]

for \(c\neq0\), even though \(0c=0\) in the underlying algebra.

Consequently, the theory must distinguish ordinary algebraic equality from
equivalence of structured division expressions.

Establishing a rigorous formal semantics for this distinction, and
determining its relationship to existing approaches such as wheels and
meadows, is the main subject for further work.


\bibliographystyle{plain}
\bibliography{references}

\end{document}
